\section{Prime-source estimate for Iwaniec's remainder weights}
\label{app:proof-aa37}
\begingroup
\setcounter{lemma}{0}
\renewcommand{\thelemma}{B.\arabic{lemma}}

This section proves the source estimate used in Section~5. Equations
(B.2)--(B.6) control the coefficients and distinctness masks.
Equations (B.7)--(B.19) treat prime detection, logarithm freezing and
boundary boxes. Equations (B.20)--(B.32) bound the retained boxes.

Fix the parameters in Theorem~\ref{thm:aa37-interior} and put
\[
 \begin{gathered}
 \tau=\sigma^9,\quad
 e_Q=(1+\sigma+\tau)\theta<\Theta(u)\le\frac{153}{256},\\
 D=N^{\theta},\quad Y=D^{\sigma^2}.
 \end{gathered}
\]
Put $s=\log D/\log z=\theta/u$, $z=N^u$, $L=\log N$,
$H=N/L^B$, and use the width $C_{\rm box}$ in
Theorem~\ref{thm:aa37-interior}. Write
$\mathcal S_n(q;a)=\mathbf1_{n\equiv a\ (q)}
-\mathbf1_{(n,q)=1}/\varphi(q)$.
All thresholds below may depend on the fixed parameters. For a product
shell $[X,2X)$, set $L_X=\log X$. The source measure is
\[
 \begin{aligned}
 \mu_D(n)&=\mathbf1_{\mathbb P}(n)
          \mathbf1_{H<n<(1-\varepsilon)N},\\
 \Delta_D(q;a)&=\sum_n\mu_D(n)
 \left(\mathbf1_{n\equiv a\ (q)}
       -\frac{\mathbf1_{(n,q)=1}}{\varphi(q)}\right).
 \end{aligned}
 \tag{B.1}
\]

\begin{lemma}[Modulus support and coefficient size]\label{lem:aa37-support}
For either sign in \textup{(W)},
\[
 \operatorname{supp}w_N^\nu\subset
 \{q<N^{e_Q}:(q,2N)=1\},\qquad
 \sup_q|w_N^\nu(q)|\le C_{\sigma,u,\theta}.
 \tag{B.2}
\]
If $[X,2X)$ meets $(H,(1-\varepsilon)N)$, then, for all sufficiently large $N$,
\[
 q<N^{e_Q}\le X^{3/5}
 \quad(q\in\operatorname{supp}w_N^\nu),\qquad H>N^{e_Q}.
 \tag{B.3}
\]
\end{lemma}

\emph{Proof.} By the definition of $\mathcal D^\nu(D)$ in \eqref{p1:eq:2.29},
every label vector in (W) satisfies the cubic label conditions of Iwaniec's
sieve remainder \cite[Theorem~4]{AA37}. Write
$s_i=\log_ND_i$ for a label vector
$\mathbf D=(D_1,\ldots,D_k)$. The cubic conditions in Iwaniec imply
$\sum_i s_i<\theta$. For the plus sign and odd $k$, use the last odd-rank condition.
For even $k$, use the condition at rank $k-1$
and $s_k\le s_{k-1}$. For the minus sign, if $k=1$ then
$s_1<\theta/2$. If $k$ is even, use the last even-rank condition.
If $k\ge3$ is odd, use rank $k-1$ and $s_k\le s_{k-1}$. Thus every
term of (W) satisfies
\[
 q<N^{\sigma \theta+(1+\tau)\sum_i s_i}
   <N^{\theta(1+\sigma+\tau)}=N^{e_Q}.
\]
All prime factors in $d$ and in the tail avoid $2N$.

Each grid coordinate has the form $D^{\sigma^2(1+\tau)^h}$ and is
less than $D^{1/2}$. Thus there are at most
$1+2\tau^{-1}\log(\sigma^{-2})$ choices per position, since
$\log(1+\tau)\ge\tau/2$. Also $k<\sigma^{-2}$ by the
last-rank cubic condition and $s_i\ge\sigma^2\theta$. With
$\tau=\sigma^9$ and $0<\sigma\le10^{-10}$, the total number of
vectors of all lengths is at most
\[
 \sum_{0\le k<\sigma^{-2}}
  \bigl(1+2\sigma^{-9}\log(\sigma^{-2})\bigr)^k
 <\exp(8\sigma^{-3}).
\]
The logarithm of the left side is at most
$\log(1+\sigma^{-2})+
\sigma^{-2}\log(1+4\sigma^{-9}\log(1/\sigma))$.
With $t=1/\sigma\ge10^{10}$, this is at most
$(3+11t^2)\log t<8t^3$, using $\log t\le t/2$.
For fixed $q$, the primes below $Y$
determine the smooth divisor $d$ uniquely. The remaining prime factors
can be assigned to the $k$ labelled positions in at most $k!$ ways.
The Iwaniec coefficients have absolute value at most $1$.
Summing over the displayed finite vector family proves the
multiplicity bound, including the empty-tail term.

If the shell meets the source interval, then $X>H/2$ and
$\log X>L-B\log L-\log2$. Impose
\[
 (3/5-e_Q)L\ge(3/5)(B\log L+\log2).
 \tag{B.4}
\]
It gives $e_QL\le(3/5)\log X$, proving the first claim
of (B.3). The second follows from $N^{1-e_Q}/L^B\to\infty$.
\hfill$\square$

Here $p>H>N^{e_Q}$, and the theorem assumes $(a,q)=1$ on the
coefficient support. So the natural comparator in (B.1) is
$\pi_{I_N}/\varphi(q)$. Thus (B.1) is the exact prime-source
discrepancy for this source.

\subsection*{Removing the distinctness masks}

Let $\widetilde w_N^\nu$ be obtained from (W) by replacing each
$K_{\mathbf D}$ by $1$, and put
$\widehat w_N^\nu(q)=\mathbf1_{(q,a)=1}\widetilde w_N^\nu(q)$.
The counting argument in Lemma~\ref{lem:aa37-support} also gives
$\sup_q|\widetilde w_N^\nu(q)|+\sup_q|\widehat w_N^\nu(q)|
\ll_{\sigma,u,\theta}1$. Its termwise product bound applies to every
unmasked term. So $\widetilde w_N^\nu$ and $\widehat w_N^\nu$ are
also supported on $\{q<N^{e_Q}:(q,2N)=1\}$, and (B.3) applies to them.
Their support need not lie in $\operatorname{supp}w_N^\nu$.
The theorem assumes $w_N^\nu(q)=0$ when $(q,a)>1$. Hence
$w_N^\nu-\widehat w_N^\nu$ is supported on the terms removed by a
distinctness mask. This holds even when an unrestricted collision
modulus is not coprime to $a$. Every such modulus is divisible by
$p^2$ for some tail prime $p\ge Y$. The coefficient difference is
$O_\sigma(1)$. Put $Q_{\rm supp}=N^{e_Q}$.
\[
 \sum_{\substack{q\le Q_{\rm supp}\\p^2\mid q\text{ for some prime }p\ge Y}}
 \left(\frac Nq+1+\frac N{\varphi(q)}\right)
 \ll \frac{N\log(2Q_{\rm supp})}{Y-1}.
 \tag{B.5}
\]
For the proof, write $q=kp^2$, use
$\varphi(kp^2)\ge\varphi(k)\varphi(p^2)$, and sum
$1/k$ and $1/\varphi(k)$ up to $Q_{\rm supp}$. Then use
$\sum_{p\ge Y}1/(p(p-1))\ll1/(Y-1)$. The left side bounds the
discrepancy of each removed term. Its progression count is at
most $N/q+1$. Consequently
\[
 \sum_q w_N^\nu(q)\Delta_D(q;a)
 =\sum_q\widehat w_N^\nu(q)\Delta_D(q;a)
   +O_\sigma(NL/Y).
 \tag{B.6}
\]
As $Y=N^{\theta\sigma^2}$, this error has a fixed power saving for each
fixed $\sigma>0$.
For a fixed Iwaniec vector, the exact factorization (4.13) survives
the projection. Replace each slot factor $\gamma_i(m)$ by
$\gamma_i(m)\mathbf1_{(m,a)=1}$. Here
$(m_1m_2m_3,a)=1$ if and only if each $(m_i,a)=1$.
The bounds $|\gamma_i|\le1$ are unchanged. Equivalently, the smooth
divisor and each tail prime are restricted to avoid $a$ within their
existing summation ranges. Thus the projected coefficient still has
the labelled form (3.4). Every modulus to which the dispersion
theorem is applied is coprime to $2a$.

\subsection*{Prime detection and the rough Heath--Brown identity}

The shells meeting the source satisfy $H/2<X<N$. Set
\[
 z_0=z_0(X)=X^{1/(\log\log X)^3},\qquad U=2X^{1/3}.
\]
For sufficiently large $N$, every $n\in[H,(1-\varepsilon)N]$ has $n\ge z_0$.
Write $\Lambda_{z_0}(n)=\Lambda(n)\mathbf1_{P^-(n)\ge z_0}$.
For every integer $n\ge z_0$, and so on the whole source range,
\[
 \mathbf1_{\mathbb P}(n)=\frac{\Lambda_{z_0}(n)}{\log n}
 -\sum_{k\ge2}\frac1k\mathbf1_{n=p^k,\ p\ge z_0}.
 \tag{B.7}
\]
The identity follows by checking primes, prime powers, and all other
integers $n\ge z_0$. It fails for primes below $z_0$, which lie
outside the source range. The proper powers contribute to the weighted discrepancy
\[
 O_\sigma(N^{1/2+o(1)}),
 \tag{B.8}
\]
in total. There are $O(N^{1/2}\log N)$ such integers in the source
range $H<n<(1-\varepsilon)N$. The theorem's range
$N\le a\le NL/3$ gives $a\ne n$ there. Each contributes
at most $O_\sigma(d(a-n)+L)$, where $d(m)$ is the divisor function.

For $n\le 8X$, the order-three rough Heath--Brown identity is
\[
 \Lambda_{z_0}(n)=\sum_{j=1}^3(-1)^{j-1}\binom3j
 \sum_{\substack{m_1\cdots m_jn_1\cdots n_j=n\\
                   m_i\le U,\ m_i,n_i\in\mathcal R_{z_0}}}
 \mu_{z_0}(m_1)\cdots\mu_{z_0}(m_j)\log n_1.
 \tag{B.9}
\]
Here $\mathcal R_{z_0}$ consists of $1$ and the $z_0$-rough integers,
and $\mu_z(m)=\mu(m)\mathbf1_{P^-(m)\ge z}$.
To verify (B.9), put
\[
 \zeta_z(s)=\sum_{r\in\mathcal R_z}r^{-s},\qquad
 M_U(s)=\sum_{m\le U}\mu_z(m)m^{-s}.
\]
The coefficients of $\zeta_zM_U$ vanish except at $1$ through degree
$U$. Thus all coefficients of $(1-\zeta_zM_U)^3$ at $n\le U^3$ vanish. Multiply the expansion of $1-(1-\zeta_zM_U)^3$ by
$-\zeta_z'/\zeta_z$ to obtain (B.9). Every variable is rough; this
includes the logarithmic variable $n_1$.

The tested indices for shell $X$ and pattern $j$ are
\[
 T_{X,j}=\left\{
 \begin{aligned}
 \mathbf v&=(m_1,\ldots,m_j,n_1,\ldots,n_j)
             \in\mathcal R_{z_0}^{2j}:\\
 &X\le\prod_i v_i<2X,\quad
   H<\prod_i v_i<(1-\varepsilon)N,\\
 &m_i\le U
 \end{aligned}\right\}.
 \tag{B.10}
\]
There are at most six variables. Their coefficient is
\[
 c_j(\mathbf v)=(-1)^{j-1}\binom3j
 \mu_{z_0}(m_1)\cdots\mu_{z_0}(m_j)\log n_1.
\]
The shell, pattern, and tuple are part of the index. The identity's
range $n\le8X$ is automatic on $T_{X,j}$.

For every nonnegative $G$ and each $j\in\{1,2,3\}$,
\[
 \sum_X\sum_{\mathbf v\in T_{X,j}}G\!\left(\prod_i v_i\right)
 \le\sum_{n<N}\tau_{2j}(n)G(n)
 \le\sum_{n<N}\tau_6(n)G(n).
 \tag{B.11}
\]
The product determines its shell. Roughness restrictions can be
dropped in this upper bound. Consequently the sum over all three
patterns is at most $3\sum_{n<N}\tau_6(n)G(n)$. The factor $3$
is absorbed in later implied constants.

For $n<N$, (B.2) and its extension to $\widehat w_N^\nu$ after (B.5)
give the pointwise majorant
\[
\left|\sum_q\widehat w_N^\nu(q)\mathcal S_n(q;a)\right|
 \le C_\sigma F_a(n),\qquad
 F_a(n)=d(a-n)+2(1+L).
 \tag*{}
\]
The congruence term counts $q\mid a-n$. The natural-comparator term
uses $\sum_{q\le N^{e_Q}}1/\varphi(q)\ll L$. On the prime source,
(B.1) equals $\sum_n\mu_D(n)\mathcal S_n(q;a)$. Apply (B.7) to this
identity while keeping $\mathcal S_n(q;a)$ attached to the tested
integer. Each HB tuple then contributes the natural PSDT discrepancy
$\mathcal S_{\prod_i v_i}(q;a)$. This holds also when its product
shares factors with $q$. The proper-power terms use the same discrepancy and
obey (B.8).

\subsection*{Equal-ratio boxes and faces}

For each product shell put
\[
 \delta_X=L_X^{-C_{\rm box}},\qquad
 J_X=\left\lceil\frac{\log2}{\log(1+\delta_X)}\right\rceil,
 \qquad \Delta_X=2^{1/J_X}-1.
\]
For sufficiently large $N$, $\delta_X\le1/3$. The equal-ratio
partition divides every dyadic coordinate interval $[2^a,2^{a+1})$
into the $J_X$ cells
\[
 [2^a(1+\Delta_X)^h,\ 2^a(1+\Delta_X)^{h+1}),
 \qquad 0\le h<J_X.
\]
It covers the full interval, including its last cell. Also
\[
 \frac{\delta_X}{3}\le\delta_X(1-2\delta_X)
 \le\Delta_X\le\delta_X.
 \tag{B.12}
\]
For the lower bound, set $\ell=\log(1+\delta_X)$. Since
$J_X<\log2/\ell+1$,
\[
 \log(1+\Delta_X)>\frac{\ell}{1+\ell/\log2}
 \ge\ell(1-\ell/\log2).
\]
Use $\delta_X-\delta_X^2/2\le\ell\le\delta_X$ and
$1/2+1/\log2<2$. This proves (B.12).

The mesh is imposed on the ambient integer axes before roughness is
tested. Keep $\{1\}$ as an atom. A product box is a Cartesian product
of at most six cells. Let $V_i$ be its lower endpoints and put
$e_i=\log_N V_i$; the atom has $V_i=1$. Define
$\lambda_{\mathcal B}=\sum_i\log V_i$. Its tested-side faces are
\[
 n=H,\quad n=(1-\varepsilon)N,\quad n=X,\quad n=2X,\quad m_i=2X^{1/3},
 \tag{B.13}
\]
where $n=\prod_i v_i$ and $m_i$ ranges over the Möbius variables.
The first two are the source cuts. The next two are product-shell
edges, and the last are the HB caps. Roughness is part of the
coefficient, not a face. A box is called boundary if one predicate in
(B.13) is not constant on it. Every other box is either empty or a
full Cartesian sum.

\begin{lemma}[Lattice thin-shell count]\label{lem:aa37-thin-shell}
Let $s\ge2$,
$a\in\mathbb Z^s\setminus\{0\}$ with $|a_i|\le2$, and
$J=\{i:a_i\ne0\}$. Let $1\le Y_0\le N$, $0<\eta_{\rm dist}\le1/4$, and let
$\mathcal L$ be $\gamma$-separated. Put
$K_\gamma=(1-e^{-\gamma/2})^{-1}$. Then
\[
\begin{split}
 \#\{\mathbf v\in\mathbb Z_{\ge1}^s:&\ \prod_i v_i\le N,
   \operatorname{dist}(\sum_i a_i\log v_i,\mathcal L)\le\eta_{\rm dist},
   \max_{i\in J}v_i\ge Y_0\}\\
 &\le s\left[4\eta_{\rm dist} K_\gamma N(1+L)^{s-1}
   +\left(\frac{4L}{\gamma}+2\right)
      \frac N{Y_0}(1+L)^{s-2}\right].
\end{split}
 \tag{B.14}
\]
For a single face, take $K_\gamma=1$ and replace
$4L/\gamma+2$ by $1$.
\end{lemma}

\emph{Proof.} Assign a tuple to the least active coordinate $i$ with
$v_i\ge Y_0$. Fix the other coordinates, with product $P$, and put
$Z=N/P$. The nonempty fibre has $P\le N/Y_0$. Let $q=|a_i|\in\{1,2\}$. The admissible values of $v_i$ lie in
intervals of relative length at most $e^{2\eta_{\rm dist}/q}-1\le4\eta_{\rm dist}$.
Their logarithmic left endpoints have spacing at least $\gamma/q$.
After clipping to $[Y_0,Z]$ and merging overlaps, there are at most
$2L/\gamma+2$ components. The separated left endpoints give total
length at most $4\eta_{\rm dist} K_\gamma Z$. Counting integer endpoints adds
at most $2L/\gamma+2$. For a singleton face there is one component.
Since $2L/\gamma+2\le4L/\gamma+2$, this gives the coefficient
$4L/\gamma+2$ in (B.14).
Sum over the fixed coordinates using
\[
 \sum_{P\le N}\frac{\tau_{s-1}(P)}P\le(1+L)^{s-1},\qquad
 \sum_{P\le N/Y_0}\tau_{s-1}(P)
 \le\frac N{Y_0}(1+L)^{s-2}.
\]
This proves (B.14). \hfill$\square$

\subsection*{Boundary and logarithm-freezing estimates}

For every shell meeting the source,
\[
 L-B\log L-\log2<\log X<L.
 \tag{B.15}
\]
Thus, since $C_{\rm box}$ is fixed, for every $\epsilon>0$,
$(\log X)^{-C_{\rm box}}\le(1+\epsilon)L^{-C_{\rm box}}$
once $N$ is large. Also $\log n$ and each nonempty box lower-corner
logarithm are $(1+o(1))L$ on the source.

A boundary box lies within $6\delta_X$ in logarithmic distance of a
face in (B.13). The source and product-shell faces give two singleton
families and the lattice $(\log2)\mathbb Z$. The cap faces give
$\log2+(\log2/3)\mathbb Z$. In each boundary tuple, an active
coordinate is at least $N^{1/7}$ for large $N$. For a product face,
the largest of at most six factors is at least
$(H/2)^{1/6}$. On a cap face the capped factor is at least
$(H/4)^{1/3}$. Apply Lemma~\ref{lem:aa37-thin-shell} with
$s\le6$, $Y_0=N^{1/7}$, and
$\eta_{\rm dist}\ll L^{-C_{\rm box}}$. The number of face families is
finite and independent of the number of boxes. Let $\mathcal U$ be the
set of boundary indices $(X,j,\mathbf v)$ across shells and patterns
for which $\mathbf v\in T_{X,j}$ lies in a boundary box. It excludes
ambient Cartesian points outside the source-restricted set (B.10).
The number of boundary indices is therefore
\[
 \#\mathcal U\ll N L^{5-C_{\rm box}}+N^{6/7}L^5
 \ll N L^{5-C_{\rm box}}.
 \tag{B.16}
\]

The elementary divisor estimates
\[
 \begin{aligned}
 \sum_{n<N}\tau_6(n)&\le N(1+L)^5,\\
 \sum_{n<N}\tau_6(n)F_a(n)^2&\ll N(1+L)^{26},\\
 \sum_{n<N}\tau_6(n)F_a(n)&\ll N(1+L)^{19}
 \end{aligned}
 \tag{B.17}
\]
follow from $\tau_6^2\le\tau_{36}$,
$d^2\le\tau_4$, $d^4\le\tau_{16}$, the divisor-sum bound
$\sum_{n\le x}\tau_k(n)\le x(1+\log x)^{k-1}$, and
Cauchy--Schwarz. For the last bound, apply Shiu's short-interval
estimate \cite{Shiu80} to $d(m)^2\le\tau_4(m)$:
$\sum_{a-N<m\le a}d(m)^2\ll N(1+L)^3$ uniformly for
$N\le a\le NL/3$. For $a\le2N$ the global divisor-sum bound
suffices. For $a>2N$ the interval length $N$ exceeds $a^{2/3}$
for all sufficiently large $N$. Combining this with
$\sum_{n<N}\tau_6(n)^2\ll N(1+L)^{35}$ gives the exponent $19$.
The bound for $F_a^2$ follows by summing $d(m)^4$ up to $a\le NL/3$.
In particular, the LOG part from
replacing
$1/\log n$ by the lower-corner value $1/\lambda_{\mathcal B}$ on a
box satisfies
\[
 \left|E_{\rm LOG}\right|\ll_\sigma NL^{18-C_{\rm box}}.
 \tag{B.18}
\]
Here $0\le\log n-\lambda_{\mathcal B}\le6\delta_X$,
$|c_j(\mathbf v)|\le3L$, and
$|1/\log n-1/\lambda_{\mathcal B}|\ll L^{-C_{\rm box}-2}$.
Sum the resulting $O(L^{-C_{\rm box}-1})F_a(n)$ over the tuples using
(B.17). A tuple belongs to exactly one box, so there is no box-count
factor.

For the boundary part, $|c_j(\mathbf v)|/\lambda_{\mathcal B}
\le3+o(1)$. Cauchy--Schwarz, (B.16), and the second estimate in (B.17)
give
\[
 |E_{\rm boundary}|\ll_\sigma NL^{16-C_{\rm box}/2}.
 \tag{B.19}
\]

\subsection*{Classification and interval lengths}

The thresholds in (B.4), together with
\[
 \sigma L\ge202(B\log L+\log2),\qquad
 \sigma L\ge202(\log2+6),
 \tag{B.20}
\]
ensure $\omega_X:=\log_N(N/X)\le\sigma/202$ and
$\kappa_X:=(\log2+6\log(1+\Delta_X))/L\le\sigma/202$. In this
appendix $\kappa_X$ plays the role of the $\kappa$ in (3.98).
For a box meeting $T_{X,j}$,
\[
 \begin{gathered}
 1-\omega_X-\kappa_X\le\sum_i e_i\le1-\omega_X+\kappa_X,\\
 e_i\le\frac{1-\omega_X}{3}+\kappa_X
 \quad\text{on each M\"obius variable}.
 \end{gathered}
 \tag{B.21}
\]
Here $\prod_i V_i\ge(\prod_i v_i)/(1+\Delta_X)^6$ and
$\prod_i V_i\le\prod_i v_i<2X$. The M\"obius bound follows from
$V_i\le m_i\le2X^{1/3}$. The corner error is independent of
the lower-cut exponent $B$.

Use the window
\[
 \mathcal W=[60\sigma,\,1/3+14\sigma].
\]
A box is Type II if some $e_i\in\mathcal W$. Select the largest such
$e_i$, with ties broken by the fixed variable order. Suppose no exponent
lies in $\mathcal W$. A box is Type I-1 when it has exactly one large
variable and that variable is not a M\"obius variable. It is Type I-2
when it has exactly two large variables and neither is a M\"obius
variable. There are no other nonboundary boxes. In fact, a large
M\"obius variable would imply
 $14\sigma\le\kappa_X-\omega_X/3\le\sigma/202$. If all variables were
small, (B.21) would give
 $1-\sigma/101\le360\sigma$. If three were large, it
would give $1+42\sigma\le1+\sigma/202$. Each is impossible.

\begin{lemma}[Lengths on retained Type-II boxes]
\label{lem:aa37-typeII-lengths}
Let $\mathcal B$
be a nonboundary Type-II box. Its selected cell is
\[
 [T,r_XT),\qquad r_X=1+\Delta_X,
\]
Let $M$ be the product of the other lower endpoints. At least one
complementary cell is nonunit, since the source product exceeds
$H>N^{1/3+15\sigma}$ eventually. There are $1\le r\le5$ such cells.
Set
\[
 I_T=(T,r_XT],\quad
 N_I=\lfloor r_XT\rfloor-\lfloor T\rfloor,\quad
 M_I=M(r_X^r-1).
 \tag{B.22}
\]
Then, for sufficiently large $N$,
\[
 N_I\asymp T L_X^{-C_{\rm box}},\qquad
 M_I\asymp (X/T)L_X^{-C_{\rm box}},
 \tag{B.23}
\]
The PSDT interval $I_T$ is contained in $(T,2T]$. It has exactly
the same selected integer coefficients as the half-open cell
$[T,r_XT)$.
\end{lemma}

\emph{Proof.} The real length of the selected cell is
$h_N=T\Delta_X$. Both
\[
 \lceil r_XT\rceil-\lceil T\rceil
 \quad\hbox{and}\quad
 N_I
\]
differ from $h_N$ by less than $1$. The only integer endpoints at
which the half-open cell and $I_T$ can differ are powers of two.
Every internal mesh endpoint is $2^a2^{h/J_X}$ with
$0<h<J_X$, and is irrational. Write $h/J_X=u/v$ in lowest terms with
$v>1$. If $2^{u/v}$ were rational, its $2$-adic valuation would give
$u=v\,v_2(2^{u/v})$, contradicting $(u,v)=1$. At powers of two the
selected rough coefficient vanishes, since $z_0>2$. Thus the two
integer conventions agree after multiplication by the coefficient.

Put
\[
 \sigma L\ge2^{130}(C_{\rm box}+1)(\log L+1).
 \tag{B.24}
\]
Then $h_N\ge T/(3L_X^{C_{\rm box}})>2$, so
$N_I\ge T/(6L_X^{C_{\rm box}})$ and
$T/N_I\le6L_X^{C_{\rm box}}\le6L^{C_{\rm box}}$. Also
$\prod_i V_i\asymp X$, hence $M\asymp X/T$ and
$M_I\asymp L_X^{-C_{\rm box}}M$, proving (B.23).
Write $x=\log_NX=1-\omega_X$, $\xi=\log_NT$, and
$\Sigma_2=\log_NS$, $\Sigma_3=\log_NR$ for the slot scales in
Lemma~\ref{lem:aa37-pack}. The scale $Q_{\rm slot}$ is the first-slot
scale in that lemma; $Q_{\rm supp}$ is the modulus cutoff in (B.5).
The inequalities in (4.9) and $M_I\asymp XT^{-1}L_X^{-C_{\rm box}}$
give
\[
\begin{aligned}
 \log_N(Q_{\rm slot}/T)
 &\le-\left(1+\frac{7271}{256}\right)\sigma<-29\sigma x,\\
 2\xi+\Sigma_2+2\Sigma_3&\le1-60\sigma,\\
 \log_N\frac{M_I X^{-89\sigma}}{Q_{\rm slot}RS}
 &\ge x-2\xi-\Sigma_2-\Sigma_3
   +\left(1+\frac{7271}{256}\right)\sigma
   -89\sigma x-o(1)\\
 &\ge \left(\frac{7271}{256}-28\right)\sigma
   -\omega_X+\Sigma_3+89\sigma\omega_X-o(1)>0.
\end{aligned}
\tag{B.25}
\]
Here $x=1-\omega_X$, $\omega_X\le\sigma/202$, and
$\log L_X/L=o(1)$ as $N\to\infty$ for fixed $\sigma,C_{\rm box}$.
The terms $o(1)$ in (B.25) are taken in this sense. For fixed $\sigma$
they are eventually smaller than $\sigma/1000$. Thus
$Q_{\rm slot}<TX^{-29\sigma}$ and
$Q_{\rm slot}RS<M_IX^{-89\sigma}$.
It follows that
\[
 N_I\ge Q_{\rm slot}X^{\sigma/3},\qquad M_I\ge X^{1/2}
 \tag{B.26}
\]
for all sufficiently large $N$. For the first inequality, the ratio
is at least
\[
 (6L_X^{C_{\rm box}})^{-1}X^{(29-1/3)\sigma};
\]
for the second, use the bound $T\le X^{1/3+15\sigma}$ of (B.30) in (B.23).

Let $m_0=M(1+r_X^r)/2$, and let $\phi_0$ be the fixed smooth cutoff in
(3.9). Then
\[
 \psi_I(m)=\phi_0\!\left(\frac{m-m_0}{M_I}\right)
\]
equals $1$ on $[M,Mr_X^r]$ and is supported on
$[M-M_I/2,Mr_X^r+M_I/2]$. It majorizes the full complementary
interval without changing its length parameter.

Since $T\ge N^{60\sigma}$ and $L_X\le L$, (B.24) gives
\[
 T^{2^{-126}}\ge12L_X^{C_{\rm box}}.
\]
Together with $N_I\ge T/(6L_X^{C_{\rm box}})$ this proves the
integer-length condition
\[
 N_I\ge(2T)^{1-2^{-126}}.
 \tag{B.27}
\]
This is the selected interval length hypothesis (LEN) in Section~3.2.
\hfill$\square$

The source cuts are constant on each retained box by the definition
of a nonboundary box. Let $n_0$ be a tested product in a retained
Type-II box, and let $T,M,r$ be as in
Lemma~\ref{lem:aa37-typeII-lengths}. Every selected integer is at most
$r_XT$. The complementary cutoff $\psi_I$ is supported below
$M[1+\tfrac32(r_X^r-1)]$. As $TM\le n_0$ and $1\le r\le5$,
their product is at most
\[
 n_0r_X[1+\tfrac32(r_X^r-1)]
 \le n_0(1+11\delta_X)
 \le(1-\varepsilon/2)N
\]
for sufficiently large $N$, using $\Delta_X\le\delta_X\to0$.
Thus on this enlarged support the upper cut leaves
\[
 a-mn\ge\varepsilon N/2\ge X(\log X)^{-C_{\rm box}}.
 \tag{B.28}
\]
The final inequality follows from $X<N$ and
$\delta_X=(\log X)^{-C_{\rm box}}<\varepsilon/2$.
Thus the resonant class is absent. Also
$a\le X^{1+\sigma/20}$ for all sufficiently large $N$.

For each fixed Iwaniec vector, Lemma~\ref{lem:aa37-pack} and the exact slot identity
(4.13) give the three-factor coefficient required by (3.5).
The smooth-only term has all three prime slots empty. Every tail factor
is at least $Y=N^{\theta\sigma^2}>N^{\ell_0\sigma^2}=u_0$.
This is the roughness threshold $u_0$ of (3.2), which bounds the tail
cutoff $u_L$ in C1. The smooth divisor is attached to the first
slot. The Section~4 packing lemma verifies the spectral inequality
(3.7) directly, as displayed in (4.9). Separately, (B.25)--(B.26)
verify the modulus and interval inequalities in (3.6); (B.27) gives
the integer-length condition (LEN), and (B.30) gives the selected-window
bounds used below. Also
\[
 \log_XN=\frac1{1-\omega_X}\le1+\frac{\sigma}{20},
\]
and $\log_XT\le1/3+15\sigma$ by the first inequality in (B.30) below.
The selected window also gives $X^{60\sigma}\le T$.

\subsection*{The interior estimate}

Let $E_{\rm int}$ be the sum of the HB contributions on retained
Type-I-1, Type-I-2, and Type-II boxes. It is taken after the finite
Iwaniec vector sum.
The arguments above prove
\begin{multline*}
 \left|\sum_qw_N^\nu(q)\Delta_D(q;a)\right|
 \\
 \le C_\sigma\bigl(
 NL/Y+N^{1/2+o(1)}+NL^{18-C_{\rm box}}\\
       +NL^{16-C_{\rm box}/2}\bigr)+|E_{\rm int}|.
 \tag{B.29}
\end{multline*}
The first two errors in (B.29) have fixed power savings. The choice
$C_{\rm box}\ge\lceil2A+32\rceil$ in Theorem~\ref{thm:aa37-interior}
makes the last two
$O(NL^{-A})$. Both estimates are over the union of tuples.

The tested measure $\mu_D$ is the prime indicator on the interval
of (B.1). It does not depend on the Iwaniec vector. The sieve factors
belong to the modulus coefficient in (W). The prime-detection and
rough Heath--Brown identities are (B.7)--(B.10). Every
variable, including the logarithmic one, is $z_0$-rough. The
prime-power terms are estimated before the box partition.

The exponent constraints (B.20)--(B.21) imply that the residual cells
are empty by the three contradictions above. The Type-II selected exponent
obeys
\[
 \begin{aligned}
 \log_XT&\le\frac{1/3+14\sigma}{1-\omega_X}
 \le1/3+15\sigma,\\
 \frac{1-\omega_X}{2}-\log_NT
 &\ge\frac16-(15+1/202)\sigma>0.
 \end{aligned}
 \tag{B.30}
\]
The second line records $T<X^{1/2}$, the Type-II upper-length condition;
the first line supplies the sharper bound used in (B.23).
By (B.2), its extension to $\widetilde w_N^\nu$ and $\widehat w_N^\nu$
after (B.5), and the parameter choice in Theorem~\ref{thm:aa37-interior},
$q<N^{e_Q}$ with $e_Q<\Theta(u)\le\Theta(0)=153/256<0.6$.
For a Type-I-1 box, (3.100) with $e_Q<153/256$ gives
$e_*-e_Q-(1-\omega_X)/10>0.3023$.
Take a Type-I-2 box. At most four of its at most six variables are
small, each with $e_i<60\sigma$. Their product $m$ has upper endpoint at
most $N^{240\sigma}(1+\Delta_X)^4\le N^{240\sigma+\kappa_X}$, so
$\log_NM\le240\sigma+\kappa_X$. Since $\log_NX=1-\omega_X$, $q<N^{e_Q}$,
$\kappa_X,\omega_X\le\sigma/202$ by (B.20), and
$0<\sigma\le10^{-10}$,
\[
 \begin{aligned}
 \log_XM&\le\frac{240\sigma+\kappa_X}{1-\omega_X}
 \le\frac{(240+1/202)\sigma}{1-\sigma/202}
 <2.401\cdot10^{-8},\\
 \log_Xq&<\frac{e_Q}{1-\omega_X}
 <\frac{153/256}{1-\sigma/202}<0.5977.
 \end{aligned}
\]
Also $N_1,N_2\ge X^{1/3}$ by the
large-variable definition, and $N_1N_2M\le2^{30}X$ by the
Heath--Brown product support and its dyadic scale convention.
Hence (3.101) holds, and with (B.30) and (3.100) the hypotheses
(3.99)--(3.101) of the Type-I lemmas hold.

The packing lemmas give, for each Iwaniec vector, the exact
three-slot coefficient and the packing inequalities (3.6)--(3.7).
They include the empty-tail term and retain repeated-prime
multiplicity before the mask error (B.6). The selected rough unit,
M\"obius and normalized-log coefficients are the three classes in
(3.8). The proof of (3.25) in Section~3 applies to each class, with
every auxiliary coprimality parameter $d$ and every subinterval.
Its bound is $T\tau(d)L^{-A_{\rm SW}}$ for every fixed
$A_{\rm SW}>0$. Its application to $\alpha^{(r)}$ and the small-norm
case in (3.47)--(3.48) give the norm-normalized BFI passage. The
complementary short-interval second moment and the remaining
dispersion steps are those of Theorem~\ref{thm:main}. On these boxes,
(B.22)--(B.28) verify their interval, smoothing and nonresonance data.
In the local theorem take $K_{\rm scale}=B+1$. In particular
$X>N/(2L^B)$ and $a\le NL/3$ imply
$a\le(2/3)XL^{B+1}\le X^{1+\sigma/20}$ for large $N$. The
comparator is the natural PSDT comparator on every HB product. No
coprimality of a composite product is asserted.

For the chosen width, the global box count in (3.91) has
$K_0=24C_{\rm box}+26$. The source labels, Iwaniec vectors and prefix
cells number $O_\sigma(1)$, and the HB coefficients are $O(1)$. There
are at most six variables. So there are $O(L^6)$ dyadic tuples, within
the $O(L^{24})$ of \S\ref{sec:pi-global}, and each dyadic tuple contains
$O(L^{6C_{\rm box}})$ short boxes, within the $O(L^{24C_{\rm box}})$ there.
This verifies (H7) for the prime source. Choose the local Type-II precision
$A'=K_0+A+5=A+24C_{\rm box}+31$ and
$A_{\rm work}=A'+4$, as in (3.92).
Then Theorem~\ref{thm:main} with (H7), that is (3.92), bounds the sum
over all Type-II boxes. The bound is $O(NL^{-A-5})$ plus a fixed
power-saving term. Lemmas~\ref{lem:typeI1}
and \ref{lem:typeI2} give $O(NL^{-A})$ for Type-I-1 and Type-I-2
when requested with $A_i=K_I+A+5$ and
$K_I=6(C_{\rm box}+1)+2$. Thus
\[
 |E_{\rm int}|\ll NL^{-A}
       +N^{1-c_\sigma}L^{K_0+C},
 \qquad c_\sigma=\ell_0\sigma^2/4>0.
 \tag{B.31}
\]
Here $C=C(\sigma,C_{\rm box},A',K_{\rm scale},\mathcal C_{\rm src})$
is the logarithmic exponent from (3.13).
Here the local spectral, restoration, endpoint and slot-localization
costs are included once in Theorem~\ref{thm:main}.

\begin{proposition}[Iwaniec prime-source estimate]\label{prop:aa37-prime-source}
Take the parameters of Theorem~\ref{thm:aa37-interior}, either Iwaniec
sign $\nu$, and a sufficiently large even $N$. Let $w_N^\nu$ be
the exact coefficient (W). Let $a$ satisfy the residue conditions of
Theorem~\ref{thm:aa37-interior}, and let $\Delta_D$ be (B.1).
First split off the LOG correction algebraically. Then every remaining
source summand is assigned exactly once to one class. The classes are
mask, proper-power, boundary, Type-II, Type-I-1 and Type-I-2.
Every retained box has the scale and coefficient properties stated
above, and
\[
 \left|\sum_qw_N^\nu(q)\Delta_D(q;a)\right|
 \ll_{\sigma,A,u,\theta,B,\varepsilon,C_{\rm box}} NL^{-A}.
 \tag{B.32}
\]
\end{proposition}

\emph{Proof.} The Iwaniec mask partition and support are (B.2)--(B.6).
The explicit identities (B.7)--(B.11) partition the remaining
tested measure into proper powers and HB tuples. The equal-ratio mesh
puts each tuple in exactly one box. A face-straddling box belongs to
the boundary class. On every other box both source predicates are constant. The
exponent-polytope argument above puts every nonempty such box in
exactly one of II, I-1 and I-2. The displayed tie rule selects
one II factor. LOG is split off algebraically before that assignment.
The errors are (B.6), (B.8), (B.18), (B.19) and (B.31). Their sum is
$O_{\sigma,A,u,\theta,B,\varepsilon,C_{\rm box}}(NL^{-A})$. \hfill$\square$

\endgroup
